\documentclass[11pt,a4paper]{article}

\usepackage[margin=1in]{geometry}
\usepackage{amsmath,amssymb}
\usepackage{booktabs}
\usepackage{hyperref}
\usepackage{xcolor}
\usepackage{graphicx}

\title{The 4D Hilbert Curve Corrected: \\
Spectral Correlation Stronger Than 3D, \\
But the Gap to Identity Remains}
\author{Rick Wesson\\
Support Intelligence, Inc.}
\date{June 2026}

\begin{document}
\maketitle

\begin{abstract}
We correct a critical error in the 4D Hilbert curve implementation used
in our earlier comparative analysis (Paper IV, Paper V).  The original
\verb|d2xyzw4D| function achieved only 50\% spatial adjacency between
consecutive curve points -- the same class of bug that affected our early
3D work.  Replacing it with the verified Butz algorithm (94.17\%
adjacency) completely reverses the previous finding: 4D Hilbert plane
decomposition now shows \emph{stronger} eigenvalue--zeta zero correlation
than 3D (Pearson $r = -0.990$ at order 4 vs $r = -0.970$ for 3D spatial).
Nevertheless, the eigenvalue magnitudes still do not converge to the zeta
zeros (MAD $\approx 0.50$, stable across orders), confirming that the gap
between correlation and identity cannot be closed by optimizing the
Hilbert construction.
\end{abstract}

\section{The Bug}

Our original \verb|d2xyzw4D| function (Paper IV, \S4) was an ad-hoc state
machine achieving exactly 50.0\% spatial adjacency between consecutive
integers along the curve.  A correct 4D Hilbert curve, like its 3D
counterpart, should have $\sim$94\% adjacency (limited only by hypercube
boundaries where the curve transitions between distant regions).

The bug was identical in nature to the 3D Hilbert curve bug discovered
and fixed earlier (Paper IV, \S3.1): both used incomplete
reflection/swap logic that failed to maintain spatial continuity across
octant boundaries.

\section{The Fix: Butz Algorithm for 4D}

We replaced the broken implementation with the verified Butz algorithm
(\S1.1 of Hamilton, ``Compact Hilbert Indices'', 2006), parameterized for
$D=4$ dimensions.  The algorithm processes 4 bits per recursion level
(one per dimension), using Gray code, rotation, and entry-point state
variables $(e, d)$ to track sub-hypercube orientation:

\begin{itemize}
\item \textbf{Gray code:} $\text{gc}(i) = i \oplus (i \gg 1)$
\item \textbf{Direction map:} $\text{dmap}(i) = g(i) \bmod 4$ where
   $g(i) = \text{trailing-zeros}(i)$
\item \textbf{Entry map:} $\text{emap}(i) = \text{gc}(2 \cdot \lfloor(i-1)/2\rfloor)$
\item \textbf{State update:} $e \leftarrow e \oplus \text{rotLeft}(\text{emap}(w), d+1)$,
   $d \leftarrow (d + \text{dmap}(w) + 1) \bmod 4$
\end{itemize}

The corrected curve achieves 94.17\% adjacency across orders 2--7,
consistent with the theoretical maximum for a 4D space-filling curve
(limited by the $1 - (1 - 2/2^n)^4$ fraction of boundary cells).

\section{Results}

\subsection{Curve-Adjacent Covariance}

\begin{table}[h]
\centering
\begin{tabular}{cccccc}
\toprule
Order & Dim & Primes & Pearson $r$ & Spearman $\rho$ & MAD \\
\midrule
4 & 16 & 6,542 & $-$0.970 & $-$1.000 & 0.518 \\
5 & 32 & 82,025 & $-$0.966 & $-$1.000 & 0.489 \\
6 & 64 & 1,077,871 & $-$0.866 & $-$1.000 & 0.518 \\
7 & 128 & 14,630,843 & $-$0.975 & $-$1.000 & 0.520 \\
\bottomrule
\end{tabular}
\caption{Corrected 4D curve-adjacent covariance: eigenvalue--zeta zero correlation.}
\end{table}

The Pearson correlation is consistently strong ($r = -0.87$ to $-0.98$)
with the exception of order 6, where a temporary dip occurs.  Order 7
recovers to $r = -0.975$, suggesting the order-6 dip is a finite-size
artifact rather than a trend.

\subsection{Spatial Adjacency Covariance}

Using 8 face-neighbors per 4D cell ($\pm x, \pm y, \pm z, \pm w$), we
construct the spatial covariance matrix analogously to the 3D case:

\begin{table}[h]
\centering
\begin{tabular}{cccccc}
\toprule
Order & Dim & Primes & Pearson $r$ & Spearman $\rho$ & MAD \\
\midrule
4 & 16 & 6,542 & $-$0.990 & $-$1.000 & 0.561 \\
5 & 32 & 82,025 & $-$0.989 & $-$1.000 & 0.537 \\
6 & 64 & 1,077,871 & $-$0.988 & $-$1.000 & 0.538 \\
7 & 128 & 14,630,843 & $-$0.961 & $-$1.000 & 0.500 \\
\bottomrule
\end{tabular}
\caption{4D spatial adjacency covariance: 8 face-neighbors per cell.}
\end{table}

Spatial adjacency achieves the strongest Pearson correlations observed:
$r = -0.990$ at order 4, staying at $r \geq -0.988$ through order 6.
This outperforms the best 3D spatial result ($r = -0.970$).  The drop to
$r = -0.961$ at order 7 may reflect the increasing sparsity of primes
(only 5.5\% of cells are prime at order 7).

\subsection{Comparison with 3D}

\begin{table}[h]
\centering
\begin{tabular}{lccc}
\toprule
Method & Orders & Best $r$ & MAD range \\
\midrule
3D curve-adjacent & 6--10 & $-$0.962 & 0.50--0.55 \\
3D spatial & 6--10 & $-$0.970 & 0.28 \\
4D curve-adjacent (corrected) & 4--7 & $-$0.975 & 0.49--0.52 \\
4D spatial & 4--7 & $-$0.990 & 0.50--0.56 \\
\bottomrule
\end{tabular}
\caption{Cross-dimensional comparison of eigenvalue--zeta zero correlation.}
\end{table}

The corrected 4D curve \emph{reverses} the previous finding that ``4D
weakens with order'' (Paper V, \S4).  That conclusion was entirely an
artifact of the broken 4D curve.  In fact, 4D spatial adjacency produces
the strongest eigenvalue--zeta zero correlation observed across any
Hilbert construction to date.

\section{Convergence Analysis}

Despite the strong correlations, the eigenvalue magnitudes do NOT
converge to the zeta zeros:

\begin{itemize}
\item \textbf{MAD stability:} The mean absolute deviation between
   normalized eigenvalue and zeta zero spectra is stable at $\approx
   0.50$--$0.56$.  Linear regression of MAD vs. order gives slope
   $-0.018$ ($p = 0.07$) for 4D spatial -- a suggestive but
   non-significant downward trend.

\item \textbf{No convergence to identity:} Even at $r = -0.990$, the
   eigenvalues differ from zeta zeros by $\sim$50\% of the spectral
   range.  This is the same ``correlation $\neq$ identity'' gap that
   we identified in the 3D case (Paper V).

\item \textbf{Pearson r varies:} The correlation fluctuates across orders
   ($-0.866$ to $-0.990$) rather than converging monotonically toward
   $-1.0$.
\end{itemize}

\section{Discussion}

\subsection{Why 4D Outperforms 3D (in Correlation)}

The 4D Hilbert curve encodes 4 bits per recursion level, mapping to 16
hyper-octants.  Primes modulo 16 occupy 8 residue classes $\{1, 3, 5, 7,
9, 11, 13, 15\}$.  The 4D Z-plane decomposition ($2^n$ planes) divides
primes more finely than 3D, potentially capturing higher-order modular
structure.

The 4D spatial approach benefits from 8 face-neighbors per cell (vs. 6
in 3D), providing richer spatial coupling information.  This may explain
the improved Pearson correlation.

\subsection{The Spearman Triviality}

The Spearman rank correlation $\rho = -1.0$ at every order, for every
method, is a \emph{trivial} consequence of eigenvalue monotonicity.  Real
symmetric matrices produce real eigenvalues; when sorted by absolute
magnitude, these are (almost always) strictly decreasing.  Zeta zeros are
strictly increasing.  Two monotonic sequences going in opposite
directions always have $\rho = -1.0$, regardless of whether they share
any deeper structure.

This is \emph{not} evidence for or against the Hilbert--P\'{o}lya
conjecture.  It is simply a property of sorted real sequences.

\subsection{What Would Close the Gap?}

The remaining obstacle is the same as in 3D (Paper V, \S2): the
eigenvalue \emph{magnitudes} must converge to the zeta zero
\emph{values}, not merely correlate with them.  For the
Hilbert--P\'{o}lya conjecture to be realized:

\begin{enumerate}
\item MAD must converge to 0 as order $\to \infty$ -- not observed.
\item The eigenvalue pair correlation function must match GUE, not
   anti-GUE -- not tested for 4D in this work but likely unchanged.
\item The operator must be proven self-adjoint in the infinite-order
   limit -- no progress.
\end{enumerate}

None of these criteria are satisfied by the corrected 4D construction.

\section{Conclusion}

The corrected 4D Hilbert curve produces eigenvalue--zeta zero correlation
of $r = -0.990$ (spatial) and $r = -0.975$ (curve-adjacent), both
\emph{stronger} than the best 3D results.  The previous finding that
``4D weakens with order'' was an artifact of a broken curve
implementation achieving only 50\% spatial adjacency.

However, the fundamental gap between correlation and identity persists.
The eigenvalues are in the correct \emph{order} but at the wrong
\emph{magnitudes} (MAD $\approx 0.50$).  No amount of optimization over
Hilbert curve constructions -- 3D, 4D, spatial, or curve-adjacent -- has
closed this gap.  The Hilbert--P\'{o}lya pathway, while tantalizingly
suggestive, cannot prove the Riemann Hypothesis with this class of
operators.

The corrected 4D results strengthen the case that the Hilbert plane
operator captures genuine spectral structure of the primes (the
correlation is too strong and too consistent to be coincidental).  But
they also reinforce the negative conclusion: capturing spectral structure
is not the same as reproducing the exact spectrum.  The true
Hilbert--P\'{o}lya operator, if it exists, must satisfy stronger
constraints than any finite-order geometric construction has achieved.

\section*{Data Availability}

All code, covariance matrices, and eigenvalue data are available in the
companion repository under \texttt{pubs/}.  Key files:

\begin{itemize}
\item \texttt{cov\_4d\_n\{4-7\}\_fixed.json} -- Corrected 4D
   curve-adjacent covariance matrices
\item \texttt{cov\_4d\_spatial\_n\{4-7\}.json} -- 4D spatial adjacency
   covariance matrices
\item \texttt{eigenvalues\_4d\_n\{4-7\}\_fixed.npy} -- Eigenvalue vectors
\end{itemize}

\end{document}
